\documentclass[10pt,a4paper,onecolumn]{IEEEtran}

\usepackage[T1]{fontenc}
\usepackage[utf8]{inputenc}
\usepackage[spanish,es-tabla]{babel}
\usepackage{lmodern,microtype,csquotes}
\usepackage{longtable,tabularx,array,booktabs,enumitem,needspace}
\usepackage[left=19mm,right=19mm,top=17mm,bottom=18mm]{geometry}
\usepackage[hidelinks]{hyperref}
\usepackage[backend=biber,style=ieee,sorting=none]{biblatex}
\addbibresource{referencias.bib}
\DefineBibliographyStrings{spanish}{bibliography={Referencias}}
\urlstyle{same}
\setlength{\parindent}{0.65em}
\setlength{\parskip}{2pt}
\setlength{\tabcolsep}{3pt}
\renewcommand{\arraystretch}{1.08}
\setlist{nosep,leftmargin=*}
\newcolumntype{Z}{>{\raggedright\arraybackslash}X}

\title{Comparación de modelos de proceso de requisitos\\
\large GA-PRAC-01 · Caso de estudio: sistema de un GAD}
\author{Cedeño Ávila Winston Damián y Paca Pilataza Luis Alberto\\
{\small Universidad Técnica Estatal de Quevedo · Facultad de Ciencias de la Computación}\\
{\small Carrera de Software · Cuarto nivel «B» · Ingeniería de Requisitos (ISR-401)}\\
{\small Grupo CPR · 7 de noviembre de 2026}}

\begin{document}
\maketitle
\begin{abstract}
Se comparan tres formas de organizar los requisitos —cascada, iterativa y ágil— mediante un caso de estudio de un GAD con contrato público. Primero se revisan siete criterios del contexto y después se identifican los actores que podrían intervenir y las actividades de ingeniería de requisitos. Es un ejercicio académico basado en supuestos; no se ha estudiado un contrato o GAD específico.
\end{abstract}
\textbf{Caso de estudio:} sistema de un GAD con contrato público. Fue elegido entre los ejemplos de la guía FORM-PC-01; todavía no consta una asignación oficial del escenario al grupo CPR.

\section{Delimitación del escenario}
La guía de clase menciona cuatro escenarios y no contamos con un dato que confirme el asignado al grupo. Por eso trabajamos con la opción del GAD como ejemplo, sin afirmar que un municipio real haya solicitado el sistema.
\textbf{Para hacer la comparación asumimos lo siguiente:}
\begin{itemize}
\item Se contrataría un sistema administrativo y existirían entregables y responsables de aceptar el trabajo.
\item No tenemos pliego ni contrato: las posibles aprobaciones y gestiones de cambios se plantean solo para comparar los modelos.
\item Podrían intervenir personal de ventanilla, una persona responsable del GAD, el área de contratación y un equipo proveedor.
\end{itemize}
Los juicios de la matriz dependen de esos supuestos y deberán revisarse si el docente proporciona otras condiciones. El proceso de requisitos se adapta al tipo de proyecto \cite{pohl2025,ireb2026}.

\section{Caracterización de los modelos}
\small
\begin{longtable}{@{}p{2.0cm}p{5.25cm}p{3.3cm}p{5.55cm}@{}}
\toprule
\textbf{Modelo} & \textbf{Fases / actividades típicas} & \textbf{Supuesto} & \textbf{Condiciones de uso}\\
\midrule\endfirsthead
\toprule
\textbf{Modelo} & \textbf{Fases / actividades típicas} & \textbf{Supuesto} & \textbf{Condiciones de uso}\\\midrule\endhead
Cascada & Levantamiento y análisis de requisitos; especificación; revisión y aprobación; desarrollo y pruebas; aceptación. & Requisitos relativamente estables y responsables capaces de aprobar una línea base. & Alcance contractual o técnico claro, cambios poco frecuentes y necesidad de documentación formal. \\
\midrule
Iterativo & Levantamiento inicial; especificar una parte; construir o simular; revisar; ajustar; repetir hasta integrar. & Se dispone de revisiones periódicas y los hallazgos pueden incorporarse por ciclos. & Partes del sistema validables por módulos, requisitos inicialmente incompletos o integración riesgosa. \\
\midrule
Ágil & Visión y backlog; refinamiento; planificación; desarrollo de incremento; revisión y adaptación del backlog. & Hay una persona que prioriza y stakeholders disponibles para retroalimentación frecuente. & Necesidades cambiantes, valor entregable por incrementos y posibilidad de revisar prioridades. \\
\midrule
\end{longtable}
\normalsize
\textit{Nota:} estas fases son una síntesis didáctica; los procesos reales pueden combinar actividades de los tres enfoques. La elicitación, documentación, validación y gestión de requisitos no desaparecen por usar ciclos o sprints \cite{ireb2026}.

\section{Criterios y matriz comparativa contextualizada}
Se emplean siete criterios para evitar elegir el modelo solo por preferencia. Cada celda explica qué ocurriría en el caso supuesto del GAD.
\textbf{Qué se evalúa en cada criterio.} 
\emph{Estabilidad y cambios}: cuánto puede variar lo solicitado y cómo se aprueban ajustes; 
\emph{stakeholders}: quién participa y qué disponibilidad tiene; 
\emph{tipo de sistema e integración}: módulos, interfaces o dependencias técnicas; 
\emph{riesgo y criticidad}: consecuencias de equivocarse; 
\emph{documentación y trazabilidad}: posibilidad de seguir cada requisito hasta su comprobación; 
\emph{validación y entrega}: cómo se comprueba lo construido; 
\emph{condiciones de aplicación}: acuerdos, tiempos y restricciones para trabajar con ese modelo.


\footnotesize
\begin{longtable}{@{}p{2.42cm}p{4.5cm}p{4.5cm}p{4.5cm}@{}}
\toprule
\textbf{Criterio} & \textbf{Cascada} & \textbf{Iterativo} & \textbf{Ágil} \\
\midrule\endfirsthead
\toprule
\textbf{Criterio} & \textbf{Cascada} & \textbf{Iterativo} & \textbf{Ágil} \\
\midrule\endhead
Estabilidad y cambios & Si el contrato fija el alcance, una línea base y solicitud formal de cambio facilitan registrar impactos en plazo y entrega. & Puede aclarar módulos con prototipos, pero cada cambio necesita cotejarse con el alcance contractual. & El backlog flexible permite ajustes frecuentes, que podrían chocar con aprobaciones y entregables pactados. \\
\midrule
Stakeholders & Actas de levantamiento y validación por responsables del GAD delimitan quién aprueba cada requisito. & Revisiones de módulos con ventanilla y sistemas permiten descubrir diferencias de interpretación. & Exige disponibilidad continua de funcionarios para ordenar backlog; esa disponibilidad no está confirmada. \\
\midrule
Tipo de sistema e integración & Permite describir trámites, perfiles y reportes antes de implementar; integración tardía es una debilidad. & Permite validar trámites por módulos sin perder acuerdos institucionales. & Incrementos ayudan a visualizar pantallas, pero deben ajustarse a la documentación exigible. \\
\midrule
Riesgo y criticidad & La aprobación de una especificación disminuye disputas de alcance, aunque errores iniciales pueden propagarse. & Pruebas por módulo detectan flujos incompletos; se requiere control formal de desvíos. & La adaptación frecuente puede afectar predictibilidad de cronograma y aceptación contractual. \\
\midrule
Documentación y trazabilidad & Matriz requisito-entregable-prueba y actas de aprobación sostienen un expediente coherente. & Exige mantener versiones de especificaciones, decisiones y pruebas en cada iteración. & Historias sin trazabilidad formal serían insuficientes para explicar cumplimiento de entregables pactados. \\
\midrule
Validación y entrega & Hitos definidos para aprobación documental y aceptación; conviene agregar prototipos tempranos. & Revisión parcial con cada dependencia del GAD antes de aceptación institucional. & Demostraciones frecuentes, siempre que la revisión tenga representación institucional autorizada. \\
\midrule
Condiciones de aplicación & Es adecuado cuando un acuerdo formal establece alcance y entregables relativamente estables. & Es más apropiado si el contrato permite aceptación parcial y revisión progresiva. & Es conveniente si hay contrato flexible, capacidad de priorización y aprobación continua. \\
\midrule
\end{longtable}
\normalsize



\section{Actores posibles y responsabilidades}
Los siguientes roles servirían para organizar el levantamiento y la revisión de requisitos. No corresponden a personas entrevistadas ni a cargos confirmados de un GAD particular.
\begin{itemize}
\item \textbf{Responsable institucional del GAD:} definir objetivos, prioridades y criterios de aceptación; autorizar decisiones según sus competencias.
\item \textbf{Personal de atención o ventanilla:} describir tareas, excepciones y necesidades de uso cotidiano.
\item \textbf{Área de tecnologías:} aclarar integraciones, permisos, seguridad, mantenimiento y restricciones técnicas.
\item \textbf{Área de contratación o asesoría jurídica:} revisar la coherencia entre requisitos, alcance acordado y procedimiento de cambios, si participa en el proyecto.
\item \textbf{Equipo proveedor:} elicitar, analizar, documentar, proponer soluciones y registrar evidencias de validación.
\item \textbf{Ciudadanía usuaria:} aportar necesidades de accesibilidad, claridad y seguimiento cuando el sistema incluya servicios dirigidos al público.
\end{itemize}

\section{Aplicación del framework de ingeniería de requisitos}
Para aplicar el framework de IR al ejemplo se separan las actividades centrales de las de gestión que acompañan todo el proceso \cite{ireb2026}. La tabla indica lo que se haría y qué documento debería quedar como resultado; ninguna actividad se presenta como ejecutada.
\small
\begin{longtable}{@{}p{3.0cm}p{6.25cm}p{5.65cm}@{}}
\toprule
\textbf{Actividad} & \textbf{Aplicación hipotética al GAD} & \textbf{Artefacto / evidencia esperada}\\
\midrule\endfirsthead
\toprule
\textbf{Actividad} & \textbf{Aplicación hipotética al GAD} & \textbf{Artefacto / evidencia esperada}\\
\midrule\endhead
Elicitación (central) & Reunir necesidades de responsables, personal operativo y tecnología; distinguir necesidades de decisiones contractuales. & Lista inicial de interesados, preguntas de entrevista, catálogo preliminar de necesidades.\\
\midrule
Documentación (central) & Redactar requisitos funcionales y de calidad con identificadores, origen y criterios verificables. & Especificación de requisitos, historias o casos de uso, glosario.\\
\midrule
Validación y negociación (central) & Revisar ambigüedades, conflictos de prioridad y criterios de aceptación con representantes autorizados. & Acta de revisión propuesta, registro de acuerdos y criterios de aceptación.\\
\midrule
Gestión de cambios (transversal) & Evaluar impacto de una solicitud sobre alcance, costo, plazo y pruebas; no asumir aprobación automática. & Solicitud de cambio, historial de versiones y decisión documentada.\\
\midrule
Trazabilidad (transversal) & Vincular requisito, origen, componente, entregable y prueba de aceptación. & Matriz de trazabilidad requisito--entregable--prueba.\\
\midrule
Priorización y seguimiento (transversal) & Ordenar necesidades por valor, riesgo y dependencia, respetando las reglas que se acuerden. & Registro de prioridades, estado y responsable de cada requisito.\\
\midrule
\end{longtable}
\normalsize
\textbf{Ejemplo ilustrativo, no requisito real:} ``El sistema permitirá consultar el estado de un trámite mediante un identificador''. Antes de aprobarlo habría que aclarar quién consulta, qué estados existen, qué datos pueden mostrarse y cómo se comprueba su cumplimiento. Este ejemplo demuestra la necesidad de elicitación, validación y trazabilidad; no describe un trámite existente.

\section{Qué nos deja esta comparación}
Para un contrato de alcance definido, cascada facilitaría aprobar una línea base, pero dejaría el riesgo de detectar errores demasiado tarde. Un modelo iterativo permitiría probar módulos antes de integrarlos, siempre que los cambios queden documentados. El enfoque ágil necesitaría funcionarios disponibles para revisar prioridades y avances; no sabemos si eso sería posible. Por tanto, la elección depende de las condiciones del proyecto, no solo de que se trate de un GAD \cite{pohl2025,iso2018}.

\section*{Uso de herramientas de apoyo}
Se utilizó ChatGPT (GPT-6) para apoyar la organización inicial del contenido, la redacción del borrador y la revisión de LaTeX. La comparación debe revisarse con la guía y las fuentes citadas; no se atribuyen entrevistas ni datos de campo.

\printbibliography
\end{document}
